% preamble-exam-zh.tex — 共享 preamble
% 用于所有 exam-zh 模板
%
% 样式策略：
%   屏幕 → 语义彩色（蓝=关键想法，红=易错，绿=路线，灰=提示/教师备注）
%   打印 → 自动降级为黑白（通过 print mode 或 monochrome 选项）

\documentclass{exam-zh}

% ---------- 基础配置 ----------
\examsetup{
  page/size = a4paper,
  page/show-head = false,
  page/show-foot = false,
  sealline/show = false,
  font = times,
  math-font = xits,
}

\usepackage{tcolorbox}
\usepackage{needspace}
\usepackage{graphicx}
\usepackage{tikz}
\usetikzlibrary{calc,intersections,angles,quotes,arrows.meta,decorations.markings,patterns.meta,backgrounds}
\usepackage{pgfplots}
\pgfplotsset{compat=1.18}
\tcbuselibrary{skins,breakable}

% ---------- 打印/屏幕颜色切换 ----------
% 用法：\documentclass[monochrome]{exam-zh} 即可切换为纯黑白
% 注意：不要使用 \if@xxx 放在 makeatother 外部；这里改成普通条件名。
\newif\ifeduprintcolor
\eduprintcolortrue

\makeatletter
\@ifclasswith{exam-zh}{monochrome}{\eduprintcolorfalse}{}
\makeatother

% ---------- 语义颜色定义 ----------
% 屏幕：有色彩区分；打印：降级为灰度
\ifeduprintcolor
  % --- 屏幕彩色模式 ---
  \definecolor{edu-blue}{RGB}{47, 82, 122}       % 关键想法
  \definecolor{edu-blue-bg}{RGB}{243, 246, 252}
  \definecolor{edu-red}{RGB}{180, 40, 40}         % 易错提醒
  \definecolor{edu-red-bg}{RGB}{255, 245, 240}
  \definecolor{edu-green}{RGB}{34, 100, 60}       % 路线图
  \definecolor{edu-green-bg}{RGB}{242, 250, 245}
  \definecolor{edu-orange}{RGB}{160, 90, 0}       % 训练目标
  \definecolor{edu-orange-bg}{RGB}{255, 248, 235}
  \definecolor{edu-gray}{RGB}{100, 100, 100}      % 提示/教师备注
  \definecolor{edu-gray-bg}{RGB}{248, 248, 248}
  \definecolor{edu-purple}{RGB}{90, 50, 120}      % 思考题
  \definecolor{edu-purple-bg}{RGB}{248, 244, 252}
\else
  % --- 打印黑白模式 ---
  \definecolor{edu-blue}{gray}{0.15}
  \definecolor{edu-blue-bg}{gray}{0.96}
  \definecolor{edu-red}{gray}{0.25}
  \definecolor{edu-red-bg}{gray}{0.97}
  \definecolor{edu-green}{gray}{0.20}
  \definecolor{edu-green-bg}{gray}{0.97}
  \definecolor{edu-orange}{gray}{0.25}
  \definecolor{edu-orange-bg}{gray}{0.98}
  \definecolor{edu-gray}{gray}{0.40}
  \definecolor{edu-gray-bg}{gray}{0.97}
  \definecolor{edu-purple}{gray}{0.30}
  \definecolor{edu-purple-bg}{gray}{0.98}
\fi
\colorlet{edu-diagram-muted}{edu-gray}

% ---------- TikZ diagram semantic macros ----------
\definecolor{edu-diagram-ink}{RGB}{17,24,39}
\definecolor{edu-diagram-marker}{RGB}{220,38,38}
\definecolor{edu-diagram-angle}{RGB}{5,150,105}
\tikzset{
  diagram point/.style={fill=edu-diagram-ink},
  diagram tick/.style={draw=edu-diagram-marker,line width=0.9pt,line cap=round},
  diagram segment/.style={draw=edu-diagram-ink,line cap=round},
  point label/.style={inner sep=1pt},
  condition label/.style={inner sep=1pt}
}
\providecommand{\DrawSegment}[3][]{\draw[diagram segment,#1] (#2) -- (#3);}
\providecommand{\DrawDashedSegment}[3][]{\draw[diagram segment,dashed,#1] (#2) -- (#3);}
\providecommand{\Triangle}[4][]{\path[#1] (#2) -- (#3) -- (#4) -- cycle;}
\providecommand{\Quadrilateral}[5][]{\path[#1] (#2) -- (#3) -- (#4) -- (#5) -- cycle;}
\providecommand{\PolygonPath}[2][]{\path[#1] #2 -- cycle;}
\providecommand{\DiagramPointRadius}{0.052cm}
\providecommand{\PointDot}[2][]{\fill[diagram point,#1] (#2) circle[radius=\DiagramPointRadius];}
\providecommand{\PointLabel}[3][]{\node[point label,#1] at (#2) {#3};}
\providecommand{\SegmentLabel}[4][]{\node[condition label,#1] at ($(#2)!0.5!(#3)$) {#4};}
\providecommand{\CoordinateTag}[3][]{\node[condition label,#1] at (#2) {#3};}
\providecommand{\AngleMark}[4][]{\pic[draw=edu-diagram-angle,line width=0.9pt,angle radius=0.48cm,#1] {angle=#2--#3--#4};}
\providecommand{\AngleLabel}[5][]{\pic["{#5}",draw=none,angle radius=0.64cm,#1] {angle=#2--#3--#4};}
\providecommand{\RightAngleMark}[4][]{\pic[draw=edu-diagram-marker,line width=0.9pt,angle radius=0.28cm,#1] {right angle=#2--#3--#4};}
\providecommand{\EqualTick}[3][]{%
  \path[postaction={decorate},decoration={markings,mark=at position .5 with {\draw[diagram tick,#1] (0pt,-4pt) -- (0pt,4pt);}}] (#2) -- (#3);%
}
\providecommand{\DoubleEqualTick}[3][]{%
  \path[postaction={decorate},decoration={markings,mark=at position .46 with {\draw[diagram tick,#1] (0pt,-4pt) -- (0pt,4pt);},mark=at position .54 with {\draw[diagram tick,#1] (0pt,-4pt) -- (0pt,4pt);}}] (#2) -- (#3);%
}
\providecommand{\TripleEqualTick}[3][]{%
  \path[postaction={decorate},decoration={markings,mark=at position .43 with {\draw[diagram tick,#1] (0pt,-4pt) -- (0pt,4pt);},mark=at position .5 with {\draw[diagram tick,#1] (0pt,-4pt) -- (0pt,4pt);},mark=at position .57 with {\draw[diagram tick,#1] (0pt,-4pt) -- (0pt,4pt);}}] (#2) -- (#3);%
}
\providecommand{\ParallelMark}[3][]{\EqualTick[#1]{#2}{#3}}
\providecommand{\NamedSegmentPath}[3]{\path[name path=#1] (#2) -- (#3);}
\providecommand{\IntersectPaths}[3]{\path[name intersections={of=#2 and #3, by=#1}];}

% ---------- 自定义教学环境 ----------
% 共性：enhanced + breakable（跨页不断裂）

% 原题卡片 — 强边框，突出显示
\newtcolorbox{problemcard}{
  enhanced, breakable,
  colback=edu-blue-bg,
  colframe=edu-blue,
  boxrule=1.2pt,
  arc=0pt,
  left=3mm, right=3mm, top=3mm, bottom=3mm,
}

% 解题路线图 — 左边线 + 浅底
\newtcolorbox{routebox}{
  enhanced, breakable,
  colback=edu-green-bg,
  colframe=edu-green!30,
  boxrule=0pt,
  borderline west={2.5pt}{0pt}{edu-green},
  left=3mm, right=3mm, top=3mm, bottom=3mm,
}

% 关键想法 — 蓝色左边线
\newtcolorbox{keyideabox}{
  enhanced, breakable,
  colback=edu-blue-bg,
  colframe=edu-blue!30,
  boxrule=0pt,
  borderline west={2.5pt}{0pt}{edu-blue},
  left=3mm, right=3mm, top=3mm, bottom=3mm,
}

% 易错提醒 — 红色左边线
\newtcolorbox{mistakebox}{
  enhanced, breakable,
  colback=edu-red-bg,
  colframe=edu-red!20,
  boxrule=0pt,
  borderline west={3pt}{0pt}{edu-red},
  left=3mm, right=3mm, top=2.5mm, bottom=2.5mm,
}

% 提示 — 灰色虚线左边线
\newtcolorbox{hintbox}{
  enhanced, breakable,
  colback=edu-gray-bg,
  colframe=edu-gray!20,
  boxrule=0pt,
  borderline west={2pt}{0pt}{edu-gray, dashed},
  left=3mm, right=3mm, top=2mm, bottom=2mm,
}

% 思考题 — 紫色虚线左边线
\newtcolorbox{thinkbox}{
  enhanced, breakable,
  colback=edu-purple-bg,
  colframe=edu-purple!20,
  boxrule=0pt,
  borderline west={2pt}{0pt}{edu-purple, dashed},
  left=2.5mm, right=2.5mm, top=2.5mm, bottom=2.5mm,
}

% 教师备注 — 灰色左边线，小字
\newtcolorbox{teachernote}{
  enhanced, breakable,
  colback=edu-gray-bg,
  colframe=edu-gray!15,
  boxrule=0pt,
  borderline west={2pt}{0pt}{edu-gray!60},
  left=2.5mm, right=2.5mm, top=2.5mm, bottom=2.5mm,
  fontupper=\small,
  colupper=black!60,
}

% 训练目标 — 橙色左边线，小字
\newtcolorbox{traininggoal}{
  enhanced, breakable,
  colback=edu-orange-bg,
  colframe=edu-orange!15,
  boxrule=0pt,
  borderline west={1.5pt}{0pt}{edu-orange},
  left=2mm, right=2mm, top=1mm, bottom=1mm,
  fontupper=\small,
}

% 答题区 — 无边框，纯留白
\newcommand{\answerarea}[1][40mm]{%
  \vspace{2mm}%
  \begin{tcolorbox}[
    enhanced, breakable,
    colback=white,
    colframe=white,
    boxrule=0pt,
    height=#1,
    valign=top,
    bottom=0mm,
    after skip=0mm,
  ]
  \end{tcolorbox}%
}

\newsavebox{\diagramtikzbox}

\newenvironment{diagramcoltikz}[2]{%
  \begin{minipage}[t]{#1}
    \vspace{0pt}\centering
    \def\diagramcaption{#2}%
    \begin{lrbox}{\diagramtikzbox}%
}{%
    \end{lrbox}%
    \usebox{\diagramtikzbox}%
    \ifx\diagramcaption\empty\else
      \par\vspace{1mm}{\scriptsize\color{edu-diagram-muted}\diagramcaption}
    \fi
  \end{minipage}%
}

\newenvironment{diagramrowtikz}[3]{%
  \begin{minipage}[t]{#1}
    \vspace{0pt}\centering
    \def\diagramlabel{#2}%
    \def\diagramcaption{#3}%
    \begin{lrbox}{\diagramtikzbox}%
}{%
    \end{lrbox}%
    \usebox{\diagramtikzbox}%
    \ifx\diagramlabel\empty\else
      \par\vspace{1mm}{\scriptsize\bfseries\diagramlabel}
    \fi
    \ifx\diagramcaption\empty\else
      \par\vspace{0.5mm}{\scriptsize\color{edu-diagram-muted}\diagramcaption}
    \fi
  \end{minipage}%
}

\newenvironment{answerareawithdiagramtikz}[3]{%
  \vspace{2mm}%
  \begin{tcolorbox}[
    enhanced, breakable,
    colback=white,
    colframe=white,
    boxrule=0pt,
    height=#1,
    valign=top,
    bottom=0mm,
    after skip=0mm,
  ]
  \begin{minipage}[t]{\dimexpr\linewidth-#2-6mm\relax}
    \vspace{0pt}
  \end{minipage}\hfill
  \begin{diagramcoltikz}{#2}{#3}
}{%
  \end{diagramcoltikz}
  \end{tcolorbox}%
}

% ---------- 字体设置 ----------
% exam-zh (ctexbook) already sets CJK fonts via ctex.
% Only override if specific fonts are needed.
% macOS: Songti SC, STHeiti, STFangsong
% Linux: SimSun, SimHei, FangSong

% ---------- 页面设置 ----------
% exam-zh (ctexbook) already loads geometry.
% Override margins if needed:
% \geometry{top=16mm, bottom=16mm, left=14mm, right=14mm}

\examsetup{
  question/show-answer = true,
  fillin/show-answer = true,
  solution/show-solution = show-stay,
}

% ---------- 教师版解答样式 ----------
\newtcolorbox{solutionblock}{
  enhanced, breakable,
  colback=edu-blue-bg,
  colframe=edu-blue!25,
  boxrule=0pt,
  borderline west={2pt}{0pt}{edu-blue!50},
  arc=0pt,
  left=3mm, right=3mm, top=2mm, bottom=2mm,
  before skip=2mm,
  after skip=3mm,
  fontupper=\small,
}

% 解答步骤编号
\newcounter{solstep}
\newcommand{\solstep}[2]{%
  \stepcounter{solstep}%
  \par\medskip
  \noindent\hangindent=6mm
  {\color{edu-blue}\bfseries\textcircled{\scriptsize\thesolstep}\enspace #1}\enspace #2\par
}

\begin{document}

\section*{矩形存在性 Fusion 综合作业}

\section*{一、模型辨析与关键量（每题 8 分，共 24 分）}


\begin{question}[points=8]
  已知 $A(0,0)$，$B(4,0)$，点 $C$ 在直线 $y=2x$ 上。比较“矩形 $ABCD$”与“菱形 $ABCD$”两种固定顺序条件，下列说法正确的是（\quad）。

  \begin{choices}
    \item 两种条件都得到 $C(4,8)$
    \item 矩形得到 $C(4,8)$，菱形得到 $C\left(\dfrac85,\dfrac{16}{5}\right)$
    \item 矩形得到 $C\left(\dfrac85,\dfrac{16}{5}\right)$，菱形得到 $C(4,8)$
    \item 两种条件都无解
  \end{choices}
  \paren[B]
  \begin{solutionblock}
    矩形固定顺序用 $BC\perp AB$，故 $x_C=4$，得 $C(4,8)$。菱形固定顺序用 $AB=BC$，设 $C(t,2t)$，列 $(t-4)^2+(2t)^2=16$，非退化解为 $t=\dfrac85$。
  \end{solutionblock}
\end{question}



\begin{question}[points=8]
  点 $P$ 在第一象限，且在反比例函数 $y=\dfrac{k}{x}$ 的图像上。以 $O(0,0)$、$P$ 以及两条坐标轴上的垂足为顶点的矩形面积为 $24$，则 $k=$ \underline{\hspace{2cm}}。

  \paren[$24$]
  \begin{solutionblock}
    第一象限中，坐标轴矩形面积等于 $xy=k$，故 $k=24$。
  \end{solutionblock}
\end{question}


\needspace{8\baselineskip}

\begin{problem}[points=8]
  一次函数 $y=x+2$ 与反比例函数 $y=\dfrac{8}{x}$ 相交。求使 $x+2\ge \dfrac{8}{x}$ 的 $x$ 的取值范围。

  \answerarea[28mm]
  \setcounter{solstep}{0}
  \begin{solutionblock}
    \textbf{答案：}$\left[-4,0\right)\cup\left[2,+\infty\right)$\par\medskip
    \solstep{求交点横坐标}{由 $x+2=\dfrac8x$ 得 $x^2+2x-8=0$，解得 $x=-4$ 或 $x=2$。}
    \solstep{分段比较}{以 $x=-4,0,2$ 分段比较函数值，得解集为 $\left[-4,0\right)\cup\left[2,+\infty\right)$。}
  \end{solutionblock}
\end{problem}


\section*{二、两模型融合：存在性 + 面积/辨析（每题 12 分，共 36 分）}

\needspace{8\baselineskip}

\begin{problem}[points=12]
  已知 $A(0,0)$，$B(6,0)$，点 $C$ 在第一象限，且在反比例函数 $y=\dfrac{k}{x}$ 的图像上。若四边形 $ABCD$ 是矩形，且矩形面积为 $18$，求 $k$ 的值以及点 $C,D$ 的坐标。

  \answerarea[48mm]
  \setcounter{solstep}{0}
  \begin{solutionblock}
    \textbf{答案：}$k=18$，$C(6,3)$，$D(0,3)$\par\medskip
    \solstep{矩形存在性入口}{$AB$ 水平且 $ABCD$ 为固定顺序矩形，所以 $BC\perp AB$，得 $x_C=6$。}
    \solstep{面积反求高度}{矩形面积为 $6y_C=18$，且 $C$ 在第一象限，故 $y_C=3$。}
    \solstep{反比例参数}{$C(6,3)$ 在 $y=\dfrac{k}{x}$ 上，所以 $k=18$。}
    \solstep{求D}{$D=A+C-B=(0,3)$。}
  \end{solutionblock}
\end{problem}


\needspace{8\baselineskip}

\begin{problem}[points=12]
  已知 $A(0,0)$，$B(4,0)$，点 $C$ 在直线 $y=2x$ 上。
\begin{enumerate}[label=(\arabic*)]
  \item 若四边形 $ABCD$ 是矩形，求 $C,D$；
  \item 若四边形 $ABCD$ 是菱形，求 $C,D$。
\end{enumerate}

  \answerarea[62mm]
  \setcounter{solstep}{0}
  \begin{solutionblock}
    \textbf{答案：}矩形：$C(4,8),D(0,8)$；菱形：$C\left(\dfrac85,\dfrac{16}{5}\right),D\left(-\dfrac{12}{5},\dfrac{16}{5}\right)$\par\medskip
    \solstep{矩形}{$AB$ 水平，固定顺序矩形得 $x_C=4$，所以 $C(4,8)$，$D=A+C-B=(0,8)$。}
    \solstep{菱形}{设 $C(t,2t)$。固定顺序菱形用 $AB=BC$，列 $(t-4)^2+(2t)^2=16$，得 $t=0$ 或 $t=\dfrac85$。舍去 $t=0$，所以 $C\left(\dfrac85,\dfrac{16}{5}\right)$。}
    \solstep{菱形第四点}{$D=A+C-B=\left(-\dfrac{12}{5},\dfrac{16}{5}\right)$。}
  \end{solutionblock}
\end{problem}


\needspace{8\baselineskip}

\begin{problem}[points=12]
  已知 $A(0,0)$，$B(6,0)$，点 $C$ 在直线 $y=x$ 上。若 $A,B,C,D$ 构成矩形，且矩形面积为 $18$，求所有可能的点 $C,D$。

  \answerarea[58mm]
  \setcounter{solstep}{0}
  \begin{solutionblock}
    \textbf{答案：}$C(3,3),D(3,-3)$\par\medskip
    \solstep{先分类}{设 $C(t,t)$。直角在 $A$ 时 $t=0$ 退化；直角在 $B$ 时 $C(6,6),D(0,6)$；直角在 $C$ 时 $C(3,3),D(3,-3)$。}
    \solstep{面积筛选}{$C(6,6)$ 对应面积 $36$，不满足；$C(3,3)$ 时相邻边长均为 $3\sqrt2$，面积为 $18$，满足。}
  \end{solutionblock}
\end{problem}


\section*{三、三模型融合与压轴（每题 10 分，共 40 分）}

\needspace{8\baselineskip}

\begin{problem}[points=10]
  已知 $A(0,0)$，$B(4,0)$，点 $C$ 在第一象限，同时在直线 $y=mx$ 与反比例函数 $y=\dfrac{k}{x}$ 的图像上。若四边形 $ABCD$ 是矩形，且矩形面积为 $20$，求 $m,k$ 的值以及点 $C,D$ 的坐标。

  \answerarea[58mm]
  \setcounter{solstep}{0}
  \begin{solutionblock}
    \textbf{答案：}$m=\dfrac54$，$k=20$，$C(4,5)$，$D(0,5)$\par\medskip
    \solstep{矩形定横坐标}{$AB$ 水平且 $ABCD$ 为矩形，所以 $x_C=4$。}
    \solstep{面积定纵坐标}{面积 $4y_C=20$，且 $C$ 在第一象限，故 $y_C=5$，所以 $C(4,5)$。}
    \solstep{回代两条函数}{$C$ 在 $y=mx$ 上，所以 $m=\dfrac54$；在 $y=\dfrac{k}{x}$ 上，所以 $k=20$。}
    \solstep{求D}{$D=A+C-B=(0,5)$。}
  \end{solutionblock}
\end{problem}


\needspace{8\baselineskip}

\begin{problem}[points=10]
  已知 $A(0,0)$，$B(4,0)$，点 $C$ 在直线 $y=2x$ 上。是否存在点 $C,D$，使四边形 $ABCD$ 同时是矩形和菱形？请说明理由。

  \answerarea[50mm]
  \setcounter{solstep}{0}
  \begin{solutionblock}
    \textbf{答案：}不存在\par\medskip
    \solstep{同时满足的含义}{固定顺序四边形同时是矩形和菱形，即为正方形，所以需 $BC\perp AB$ 且 $BC=AB=4$。}
    \solstep{矩形候选}{$BC\perp AB$ 得 $x_C=4$。因 $C$ 在 $y=2x$ 上，$C(4,8)$，此时 $BC=8\ne4$。}
    \solstep{结论}{因此不存在满足条件的 $C,D$。}
  \end{solutionblock}
\end{problem}


\needspace{8\baselineskip}

\begin{problem}[points=10]
  一次函数 $y=x+2$ 与反比例函数 $y=\dfrac{8}{x}$ 交于两点。
\begin{enumerate}[label=(\arabic*)]
  \item 求两个交点；
  \item 取第一象限交点 $P$，以 $O(0,0)$、$P$ 及两条坐标轴上的垂足为顶点作矩形，求矩形面积；
  \item 解不等式 $x+2\ge\dfrac{8}{x}$。
\end{enumerate}

  \answerarea[70mm]
  \setcounter{solstep}{0}
  \begin{solutionblock}
    \textbf{答案：}交点 $(2,4)$、$(-4,-2)$；面积 $8$；不等式解集 $\left[-4,0\right)\cup\left[2,+\infty\right)$\par\medskip
    \solstep{求交点}{由 $x+2=\dfrac8x$，得 $x^2+2x-8=0$，解得 $x=2$ 或 $x=-4$。交点为 $(2,4)$、$(-4,-2)$。}
    \solstep{矩形面积}{第一象限交点 $P(2,4)$，坐标轴矩形面积为 $2\times4=8$。}
    \solstep{解不等式}{以 $x=-4,0,2$ 分段比较，得 $x+2\ge\dfrac8x$ 的解集为 $\left[-4,0\right)\cup\left[2,+\infty\right)$。}
  \end{solutionblock}
\end{problem}


\needspace{8\baselineskip}

\begin{problem}[points=10]
  已知 $A(0,0)$，$B(6,0)$。
\begin{enumerate}[label=(\arabic*)]
  \item 点 $C$ 在第一象限的反比例函数 $y=\dfrac{k}{x}$ 上，四边形 $ABCD$ 是矩形，且面积为 $24$。求 $k,C,D$；
  \item 点 $E$ 在直线 $y=2x$ 上，四边形 $ABEF$ 是菱形。求 $E,F$；
  \item 比较两个四边形的面积。
\end{enumerate}

  \answerarea[75mm]
  \setcounter{solstep}{0}
  \begin{solutionblock}
    \textbf{答案：}$k=24$，$C(6,4)$，$D(0,4)$；$E\left(\dfrac{12}{5},\dfrac{24}{5}\right)$，$F\left(-\dfrac{18}{5},\dfrac{24}{5}\right)$；菱形面积 $\dfrac{144}{5}$ 大于矩形面积 $24$\par\medskip
    \solstep{矩形部分}{$ABCD$ 为固定顺序矩形，得 $x_C=6$。面积 $6y_C=24$，故 $C(6,4)$，$D(0,4)$，且 $k=24$。}
    \solstep{菱形部分}{设 $E(t,2t)$。菱形 $ABEF$ 中 $AB=BE$，列 $(t-6)^2+(2t)^2=36$，得 $t=0$ 或 $t=\dfrac{12}{5}$。舍去退化解，得 $E\left(\dfrac{12}{5},\dfrac{24}{5}\right)$。}
    \solstep{求F和比较面积}{$F=A+E-B=\left(-\dfrac{18}{5},\dfrac{24}{5}\right)$。矩形面积为 $24$；菱形以 $AB$ 为底，高为 $\dfrac{24}{5}$，面积为 $6\cdot\dfrac{24}{5}=\dfrac{144}{5}$，所以菱形面积更大。}
  \end{solutionblock}
\end{problem}


\clearpage
\section*{答案速查}




\end{document}
